\documentclass[12pt, a4paper]{article}

\usepackage[utf8]{inputenc}
\usepackage[T1]{fontenc}
\usepackage[italian]{babel}
\usepackage{geometry}
\usepackage{hyperref}
\usepackage{graphicx}
\usepackage{booktabs}
\usepackage{longtable}
\usepackage{array}
\usepackage{xcolor}
\usepackage{listings}
\usepackage{enumitem}
\usepackage{titlesec}
\usepackage{fancyhdr}
\usepackage{setspace}
\usepackage{microtype}
\usepackage{parskip}
\usepackage{float}
\usepackage{caption}
\usepackage{subcaption}
\usepackage{mdframed}
\usepackage{lmodern}
\usepackage{amsmath}

\geometry{
    top=2.5cm,
    bottom=2.5cm,
    left=3cm,
    right=2.5cm
}

\setstretch{1.4}

\definecolor{accent}{HTML}{185FA5}
\definecolor{codebg}{HTML}{F8F9FA}
\definecolor{codeframe}{HTML}{DEE2E6}
\definecolor{green}{HTML}{166534}
\definecolor{red}{HTML}{991B1B}
\definecolor{gray}{HTML}{6B7280}

\hypersetup{
    colorlinks = true,
    linkcolor  = accent,
    urlcolor   = accent,
    citecolor  = accent,
    pdftitle   = {Relazione ParkShare},
    pdfauthor  = {[NOME STUDENTE]},
}

\lstdefinestyle{codestyle}{
    backgroundcolor  = \color{codebg},
    frame            = single,
    rulecolor        = \color{codeframe},
    basicstyle       = \ttfamily\small,
    keywordstyle     = \color{accent}\bfseries,
    commentstyle     = \color{gray}\itshape,
    stringstyle      = \color{green},
    numbers          = left,
    numberstyle      = \tiny\color{gray},
    numbersep        = 8pt,
    breaklines       = true,
    breakatwhitespace= true,
    tabsize          = 4,
    showstringspaces = false,
    captionpos       = b,
}
\lstset{style=codestyle}

\lstdefinelanguage{JavaScript}{
  keywords={break, case, catch, continue, debugger, default, delete,
    do, else, false, finally, for, function, if, in, instanceof,
    new, null, return, switch, this, throw, true, try, typeof,
    var, void, while, with, let, const, async, await, class,
    import, export, from, of},
  keywordstyle=\color{accent}\bfseries,
  ndkeywords={},
  sensitive=true,
  comment=[l]{//},
  morecomment=[s]{/*}{*/},
  morestring=[b]',
  morestring=[b]",
  morestring=[b]`
}

\pagestyle{fancy}
\fancyhf{}
\fancyhead[L]{\small\textcolor{gray}{ParkShare — Relazione di Progetto}}
\fancyhead[R]{\small\textcolor{gray}{\leftmark}}
\fancyfoot[C]{\thepage}
\renewcommand{\headrulewidth}{0.4pt}

\titleformat{\section}
  {\large\bfseries\color{accent}}{\thesection}{1em}{}[\titlerule]
\titleformat{\subsection}
  {\normalsize\bfseries\color{accent!80!black}}{\thesubsection}{1em}{}
\titleformat{\subsubsection}
  {\normalsize\bfseries\color{gray}}{\thesubsubsection}{1em}{}

\newmdenv[
    linecolor=accent,
    backgroundcolor=accent!5,
    linewidth=1.5pt,
    leftline=true,
    rightline=false,
    topline=false,
    bottomline=false,
    innerleftmargin=10pt,
    innerrightmargin=10pt,
    innertopmargin=8pt,
    innerbottommargin=8pt,
]{infobox}

\begin{document}

\begin{titlepage}
\centering

\vspace*{2cm}

{\LARGE\bfseries\color{accent} Università degli Studi\par}
{\large Dipartimento di Informatica\par}
{\large Corso di Programmazione Web\par}

\vspace{2cm}

\rule{\linewidth}{1.5pt}\par
\vspace{0.5cm}
{\Huge\bfseries ParkShare\par}
\vspace{0.3cm}
{\Large\color{gray} Una piattaforma web open-source\\ per la condivisione di parcheggi privati\par}
\vspace{0.5cm}
\rule{\linewidth}{1.5pt}\par

\vspace{2.5cm}

\begin{tabular}{rl}
    \textbf{Studente:}    & Sandi Russo \\[4pt]
    \textbf{Matricola:}   & 553675 \\[4pt]
    \textbf{Docente:}     & Andrea Nucita \\[4pt]
    \textbf{A.A.:}        & 2025/2026 \\
\end{tabular}

\end{titlepage}

\newpage
\tableofcontents
\newpage

\section{Introduzione}

\subsection{Contesto e motivazione}

Nelle città moderne, il parcheggio rappresenta uno dei problemi logistici più ricorrenti per i
guidatori: la ricerca di un posto libero genera traffico, perdita di tempo e frustrazione, con
effetti negativi sia sull'esperienza dell'utente sia sull'ambiente. Parallelamente, molti
proprietari di abitazioni o attività commerciali dispongono di posti auto privati che rimangono
inutilizzati per la maggior parte della giornata.

L'idea alla base di \textbf{ParkShare} nasce proprio da questa asimmetria: collegare chi cerca
un parcheggio con chi ne possiede uno disponibile, creando un \emph{marketplace} di parcheggi
privati accessibile via browser, senza installazione di applicazioni, senza costi di accesso per
gli utenti e senza dipendenza da servizi commerciali a pagamento.

\subsection{Obiettivi del progetto}

Il progetto si propone di raggiungere i seguenti obiettivi:

\begin{enumerate}[leftmargin=*, label=\textbf{\arabic*.}]
    \item \textbf{Funzionalità complete end-to-end}: registrazione utente, pubblicazione
    annunci, prenotazione con validazione oraria, sistema di recensioni, generazione di
    ricevute PDF.

    \item \textbf{Indipendenza da API proprietarie}: tutta la componente cartografica è
    basata su tecnologie open-source (\textbf{Leaflet.js} + \textbf{OpenStreetMap}), senza
    richiedere chiavi API, senza costi di utilizzo e senza inviare dati degli utenti a
    terze parti.

    \item \textbf{Accessibilità universale}: l'interfaccia supporta modalità chiara e scura
    (con persistenza della scelta), è completamente responsiva per dispositivi mobili,
    rispetta le linee guida WCAG per l'accessibilità (attributi ARIA, navigazione da
    tastiera) e onora la preferenza \texttt{prefers-reduced-motion} per gli utenti con
    sensibilità alle animazioni.

    \item \textbf{Sicurezza applicativa}: gestione sicura delle credenziali tramite file
    \texttt{.env}, password cifrate con \texttt{bcrypt}, query parametrizzate contro
    SQL injection, rate limiting sul login.

    \item \textbf{Qualità del codice}: architettura a strati con pattern Repository,
    separazione netta tra logica di business, accesso ai dati e presentazione.
\end{enumerate}


\section{Stato dell'Arte}
\label{sec:sota}

\subsection{Il problema del parcheggio urbano}

Il problema della sosta urbana è ampiamente documentato in letteratura. Studi condotti in
diverse città europee stimano che una quota compresa tra il 20\% e il 30\% del traffico
cittadino sia generato da veicoli alla ricerca di parcheggio. In Italia, la
situazione è aggravata dall'elevata densità veicolare nei centri storici, dalla scarsità di
parcheggi pubblici e dalla presenza di numerosi posti auto privati non sfruttati in orari di
punta.

\subsection{Soluzioni commerciali esistenti}

Negli ultimi anni sono emersi diversi servizi digitali per la gestione della sosta. I più noti
a livello internazionale sono:

\begin{description}[leftmargin=0pt]
    \item[\textbf{JustPark} (UK)] Piattaforma peer-to-peer per affitto di posti auto privati.
    Utilizza Google Maps API per la visualizzazione cartografica e addebita una commissione
    del 15\% su ogni transazione.

    \item[\textbf{SpotHero / ParkWhiz} (USA)] Aggregatori di parcheggi prevalentemente
    commerciali. Richiedono integrazione con sistemi di pagamento proprietari e sono
    geograficamente limitati al mercato nordamericano.

    \item[\textbf{EasyPark}] Applicazione diffusa anche in Italia per il pagamento della
    sosta su strisce blu. Non gestisce parcheggi privati e richiede un account a pagamento
    per i gestori.

    \item[\textbf{ParkU}] Piattaforma italiana di parking sharing, acquisita nel 2016
    da EasyPark. Non più disponibile come servizio indipendente.
\end{description}

\subsubsection{Limitazioni delle soluzioni esistenti}

Tutte le soluzioni commerciali sopra citate presentano una o più delle seguenti criticità:

\begin{itemize}
    \item \textbf{Dipendenza da API proprietarie a pagamento}: la quasi totalità utilizza
    Google Maps Platform, il cui piano gratuito è limitato a \$200/mese di credito.

    \item \textbf{Trasmissione di dati a terze parti}: l'utilizzo di Google Maps o Here Maps
    implica che le coordinate geografiche degli utenti vengano elaborate da server di
    proprietà di grandi aziende tecnologiche, con implicazioni per la privacy non sempre
    accettabili.

    \item \textbf{Lock-in tecnologico}: costruire un'applicazione su API commerciali
    significa legarsi a politiche di pricing che possono cambiare unilateralmente, come
    dimostrato dall'aumento dei prezzi di Google Maps nel 2018, che costrinse molti
    sviluppatori a migrare con urgenza.

    \item \textbf{Accessibilità limitata}: poche piattaforme supportano nativamente
    la modalità scura, la navigazione da tastiera o rispettano le preferenze di sistema
    per la riduzione del movimento.
\end{itemize}

\subsection{Tecnologie cartografiche open-source}

Di fronte alle limitazioni sopra descritte, il panorama open-source offre alternative
mature e affidabili.

\textbf{OpenStreetMap} (OSM) è un progetto collaborativo avviato nel 2004 che mette a
disposizione dati geografici liberamente utilizzabili sotto licenza ODbL. Con oltre
10 milioni di utenti contributori nel mondo, rappresenta il dataset geografico aperto
più completo disponibile.

\textbf{Leaflet.js} è una libreria JavaScript open-source (licenza BSD-2) specializzata
nella visualizzazione di mappe interattive. Con soli \textasciitilde 42~KB è significativamente più leggera delle controparti commerciali (Google Maps JS API supera
i 300~KB), garantisce prestazioni elevate su dispositivi mobili e non richiede
alcuna chiave API.

\textbf{Nominatim} è il servizio di geocoding (indirizzo $\to$ coordinate e viceversa)
ufficiale di OpenStreetMap, utilizzato in ParkShare per la funzione di rilevamento
automatico dell'indirizzo dal segnale GPS. Come Leaflet, è completamente gratuito
e non trasmette dati identificativi a servizi terzi.

\subsection{Servizi web: REST vs SOAP}

La scelta tra le due principali architetture per i servizi web è caduta su
\textbf{REST} (Representational State Transfer) per le seguenti ragioni:

\begin{itemize}
    \item \textbf{Semplicità}: REST utilizza i metodi HTTP standard
    (\texttt{GET}, \texttt{POST}, \texttt{PUT}, \texttt{DELETE}) in modo semanticamente
    coerente con l'operazione da compiere, senza richiedere un livello di overhead
    aggiuntivo come l'envelope SOAP.

    \item \textbf{Formato dei dati}: JSON è notevolmente più compatto di XML e nativamene
    supportato da JavaScript, eliminando la necessità di parser dedicati.
\end{itemize}

\section{Materiali e Metodi}
\label{sec:metodi}

\subsection{Architettura del sistema}

ParkShare segue un'architettura \textbf{client-server a tre livelli} (three-tier):

\begin{enumerate}
    \item \textbf{Presentation tier} (client): browser web con HTML5, CSS3 e JavaScript
    vanilla, responsabile del rendering dell'interfaccia e della comunicazione con il tier
    successivo tramite chiamate \texttt{fetch} asincrone.

    \item \textbf{Application tier} (server): PHP 8 che espone endpoint RESTful, esegue
    la logica di business, valida gli input e gestisce la sessione utente.

    \item \textbf{Data tier}: database MySQL con schema relazionale, accessibile
    esclusivamente tramite lo strato applicativo.
\end{enumerate}

\begin{figure}[H]
\centering
\begin{mdframed}[linecolor=codeframe, linewidth=1pt]
\begin{center}
\texttt{[Browser]} $\xrightarrow{\text{HTTP/JSON}}$ \texttt{[PHP API]} $\xrightarrow{\text{SQL}}$ \texttt{[MySQL]}
\end{center}
\vspace{0.3cm}
\small
\textbf{Client}: HTML, CSS, JS (Leaflet, fetch API)\\
\textbf{Server}: PHP 8, OOP, pattern Repository\\
\textbf{Database}: MySQL, prepared statements
\end{mdframed}
\caption{Schema dell'architettura a tre livelli di ParkShare.}
\end{figure}

La struttura delle directory riflette questa separazione:

\begin{lstlisting}[language=bash, caption={Struttura del progetto ParkShare.}]
PROGRAMMAZIONE_WEB/
|- api/
|- classes/
|- config/
|- includes/
|- pages/
|- assets/
|  |- css/
|  \- js/
|- .env
\- .env.example
\end{lstlisting}

\subsection{Scelte tecnologiche lato client}

\subsubsection{JavaScript Vanilla senza framework}

Una scelta progettuale deliberata è stata quella di \textbf{non adottare framework
JavaScript} come React, Vue o Angular. Le motivazioni sono molteplici:

\begin{itemize}
    \item \textbf{Comprensione dei fondamentali}: un progetto universitario ha lo scopo
    primario di dimostrare la padronanza delle tecnologie di base. Delegare la gestione
    del DOM, degli eventi e delle chiamate asincrone a un framework nasconde i meccanismi
    sottostanti.

    \item \textbf{Zero dipendenze da npm}: nessun \texttt{node\_modules}, nessun
    bundler (webpack, Vite), nessun processo di build. Il progetto si esegue servendo
    i file statici da un qualsiasi server PHP, senza configurazioni aggiuntive.

    \item \textbf{Longevità}: il codice scritto con API Web native (Fetch API, DOM API,
    ES2020+) non diventa obsoleto con i cicli di rilascio dei framework.
\end{itemize}

\subsubsection{Gestione asincrona con Fetch API}

Tutta la comunicazione con il server avviene tramite la \textbf{Fetch API} con la
sintassi \texttt{async/await}, che rende il flusso asincrono leggibile quanto del
codice sincrono. Le funzioni helper \texttt{postJSON}, \texttt{putJSON} e
\texttt{deleteJSON} centralizzano la gestione degli errori di rete:

\begin{lstlisting}[language=JavaScript, caption={Funzione helper per richieste POST in JSON.}]
async function postJSON(url, body) {
    try {
        const res = await fetch(url, {
            method: 'POST',
            headers: { 'Content-Type': 'application/json' },
            body: JSON.stringify(body)
        });
        return await res.json();
    } catch {
        return { success: false,
                 error: 'Errore di rete. Controlla la connessione.' };
    }
}
\end{lstlisting}

\subsection{Scelte tecnologiche lato server}

\subsubsection{PHP senza framework}

Analogamente al lato client, si è scelto di non adottare framework PHP come Laravel o
Symfony. Questo permette di:

\begin{itemize}
    \item Dimostrare la comprensione diretta di PHP: routing manuale, gestione delle
    sessioni, preparazione delle query, upload di file.
    \item Mantenere un'installazione minimale: l'unica dipendenza esterna è
    \textbf{dompdf} (via Composer) per la generazione dei PDF.
    \item Controllare ogni aspetto della sicurezza senza fare affidamento
    ``cieco'' sulle protezioni automatiche di un framework.
\end{itemize}

\subsubsection{Pattern Repository}

La logica di accesso al database è incapsulata nel pattern \textbf{Repository},
implementato tramite una classe astratta base:

\begin{lstlisting}[language=PHP, caption={Classe Repository — metodi di accesso generici al DB.}]
abstract class Repository {
    protected static function fetchOne(
        string $sql, string $types, mixed ...$params
    ): ?array { ... }

    protected static function fetchAll(
        string $sql, string $types, mixed ...$params
    ): array { ... }
}
\end{lstlisting}

Le classi di dominio (\texttt{User}, \texttt{Parcheggio}, \texttt{Prenotazione},
\texttt{Recensione}) estendono \texttt{Repository} e implementano la logica
specifica. Questo garantisce che ogni query al database passi obbligatoriamente
per i prepared statements, eliminando alla radice il rischio di SQL injection.

\subsubsection{Singleton per la connessione al database}

La classe \texttt{Database} implementa il pattern \textbf{Singleton} per garantire
che esista una sola connessione MySQLi per richiesta HTTP, evitando lo spreco di
risorse derivante dall'apertura di connessioni multiple:

\begin{lstlisting}[language=PHP, caption={Pattern Singleton per la connessione al database.}]
class Database {
    private static ?Database $instance = null;
    private mysqli $connection;

    private function __construct() {
        $this->connection = new mysqli(
            DB_HOST, DB_USER, DB_PASS, DB_NAME
        );
    }

    public static function getInstance(): Database {
        if (self::$instance === null) {
            self::$instance = new Database();
        }
        return self::$instance;
    }
}
\end{lstlisting}

\subsection{Il servizio web RESTful}

\subsubsection{Principi di progettazione}

L'API di ParkShare rispetta i principi REST di Roy Fielding: risorse identificate da
URI, operazioni espresse tramite metodi HTTP semantici, comunicazione stateless,
risposte in formato JSON.

\begin{table}[H]
\centering
\caption{Endpoint principali dell'API RESTful di ParkShare.}
\renewcommand{\arraystretch}{1.3}
\begin{tabular}{llll}
\toprule
\textbf{Metodo} & \textbf{Endpoint} & \textbf{Azione} & \textbf{Auth?} \\
\midrule
POST & \texttt{/api/auth.php} (login)     & Autenticazione utente      & No  \\
POST & \texttt{/api/auth.php} (register)  & Registrazione utente       & No  \\
POST & \texttt{/api/auth.php} (logout)    & Terminazione sessione       & Sì  \\
GET  & \texttt{/api/parcheggi.php?id=N}   & Dettaglio parcheggio       & No  \\
GET  & \texttt{/api/parcheggi.php?pagina=N}& Lista paginata             & No  \\
GET  & \texttt{/api/parcheggi.php?lat=\&lon=}& Ricerca per posizione   & No  \\
POST & \texttt{/api/parcheggi.php}        & Crea parcheggio            & Sì  \\
PUT  & \texttt{/api/parcheggi.php}        & Modifica parcheggio        & Sì  \\
DELETE & \texttt{/api/parcheggi.php}      & Elimina parcheggio         & Sì  \\
POST & \texttt{/api/prenotazioni.php}     & Crea prenotazione          & Sì  \\
PUT  & \texttt{/api/prenotazioni.php}     & Aggiorna stato prenotazione& Sì  \\
GET  & \texttt{/api/recensioni.php}       & Recensioni di un parcheggio& No  \\
POST & \texttt{/api/recensioni.php}       & Crea recensione            & Sì  \\
DELETE & \texttt{/api/recensioni.php}     & Elimina recensione         & Sì  \\
POST & \texttt{/api/utenti.php}           & Aggiorna profilo utente    & Sì  \\
GET  & \texttt{/api/ricevuta.php?id=N}    & Genera ricevuta PDF        & Sì  \\
\bottomrule
\end{tabular}
\end{table}

\subsubsection{Formato delle risposte}

Tutte le risposte adottano una struttura JSON uniforme:

\begin{lstlisting}[language=bash, caption={Struttura delle risposte API — successo ed errore.}]
// Risposta di successo
{ "success": true, "data": { ... } }

// Risposta di errore
{ "error": "Messaggio di errore leggibile dall'utente" }

// Risposta paginata (endpoint /api/parcheggi.php?pagina=N)
{
  "success": true,
  "data": [ ... ],
  "totale": 45,
  "pagina": 2,
  "perPagina": 12,
  "hasMore": true
}
\end{lstlisting}

\subsection{Database relazionale}

\subsubsection{Schema Entità-Relazione}

Il database è composto da cinque tabelle principali, con relazioni che garantiscono
l'integrità referenziale:

\begin{itemize}
    \item \textbf{UTENTE}: \texttt{ID\_UTENTE}, \texttt{NOME}, \texttt{COGNOME},
    \texttt{EMAIL} (unique), \texttt{PASSWORD\_HASH}.

    \item \textbf{PARCHEGGIO}: \texttt{ID\_PARCHEGGIO}, \texttt{ID\_PROPRIETARIO}
    (FK $\to$ UTENTE), \texttt{TITOLO}, \texttt{DESCRIZIONE}, \texttt{INDIRIZZO},
    \texttt{LAT}, \texttt{LON}, \texttt{COSTO\_ORARIO}, \texttt{ORA\_APERTURA},
    \texttt{ORA\_CHIUSURA}, \texttt{FOTO}.

    \item \textbf{DISPONIBILITA}: \texttt{ID\_PARCHEGGIO} (FK), \texttt{GIORNO}
    (0=Lun \ldots 6=Dom). Relazione N:M tra parcheggio e giorni della settimana.

    \item \textbf{PRENOTAZIONE}: \texttt{ID\_PRENOTAZIONE}, \texttt{ID\_UTENTE} (FK),
    \texttt{ID\_PARCHEGGIO} (FK), \texttt{DATA}, \texttt{ORA\_INIZIO},
    \texttt{ORA\_FINE}, \texttt{IMPORTO\_TOTALE}, \texttt{STATO}
    (\texttt{in\_attesa}~$|$~\texttt{confermata}~$|$~\texttt{cancellata}~$|$~\texttt{completata}).

    \item \textbf{RECENSIONE}: \texttt{ID\_RECENSIONE}, \texttt{ID\_UTENTE} (FK),
    \texttt{ID\_PARCHEGGIO} (FK), \texttt{ID\_PRENOTAZIONE} (FK, unique — una
    sola recensione per prenotazione), \texttt{VOTO} (1--5), \texttt{COMMENTO},
    \texttt{DATA}.
\end{itemize}

\subsubsection{Ricerca geografica con la formula di Haversine}

La funzionalità ``cerca vicino a me'' utilizza la \textbf{formula di Haversine}
direttamente in una query MySQL per calcolare la distanza in chilometri tra le
coordinate GPS dell'utente e quelle di ogni parcheggio:

\begin{lstlisting}[language=SQL, caption={Ricerca geografica con Haversine e bounding box.}]
SELECT *,
  (6371 * ACOS(
    COS(RADIANS(?)) * COS(RADIANS(LAT)) *
    COS(RADIANS(LON) - RADIANS(?)) +
    SIN(RADIANS(?)) * SIN(RADIANS(LAT))
  )) AS distanza
FROM PARCHEGGIO
WHERE LAT BETWEEN ? AND ?
  AND LON BETWEEN ? AND ?
HAVING distanza <= ?
ORDER BY distanza ASC
\end{lstlisting}

Il bounding box approssimativo riduce il numero di righe su cui calcolare la funzione
trigonometrica, migliorando le prestazioni al crescere del numero di parcheggi.

\subsection{Componente cartografica: Leaflet.js e OpenStreetMap}

\subsubsection{Motivazioni della scelta}

La decisione di non utilizzare Google Maps API e di affidarsi invece a
\textbf{Leaflet.js} con tile di \textbf{OpenStreetMap} è stata motivata da tre
ragioni fondamentali:

\begin{enumerate}
    \item \textbf{Costo zero e nessuna chiave API}: il progetto funziona
    immediatamente senza registrazioni, senza limiti di utilizzo mensile e senza
    rischi di blocco per superamento della quota gratuita.

    \item \textbf{Privacy degli utenti}: le richieste di tile vengono inviate ai
    server OSM (o a mirror della community come \texttt{tile.openstreetmap.org}),
    non ai server di Google. Nessun dato di navigazione viene trasmesso a terze
    parti commerciali.

    \item \textbf{Leggerezza}: Leaflet pesa $\sim$42~KB gzippati contro i
    $>$300~KB di Google Maps JS API v3, con impatto diretto sui tempi di
    caricamento su connessioni lente.
\end{enumerate}

\subsubsection{Reverse geocoding con Nominatim}

Per la compilazione automatica dell'indirizzo durante l'aggiunta di un parcheggio,
si utilizza il servizio \textbf{Nominatim} di OpenStreetMap, che converte le
coordinate GPS in un indirizzo leggibile:

\begin{lstlisting}[language=JavaScript, caption={Chiamata a Nominatim per il reverse geocoding.}]
const res = await fetch(
    `https://nominatim.openstreetmap.org/reverse` +
    `?lat=${lat}&lon=${lon}&format=json`,
    { headers: { 'Accept-Language': 'it' } }
);
const dati = await res.json();
const indirizzo = `${dati.address.road}, ${dati.address.city}`;
\end{lstlisting}

\subsection{Accessibilità e usabilità}

\subsubsection{Tema chiaro e scuro con persistenza}

ParkShare supporta nativamente tre modalità di visualizzazione:

\begin{enumerate}
    \item \textbf{Automatica}: rileva la preferenza di sistema tramite la media
    query CSS \texttt{prefers-color-scheme} e imposta il tema corrispondente.
    \item \textbf{Scura manuale}: l'utente attiva il tema scuro tramite il bottone in navbar; la scelta viene salvata in \texttt{localStorage} e persiste
    tra sessioni diverse.
    \item \textbf{Chiara manuale}: forza il tema chiaro anche su sistemi con
    preferenza di sistema impostata su scuro.
\end{enumerate}

L'implementazione sfrutta variabili CSS custom (\texttt{--bg}, \texttt{--surface},
\texttt{--text}, ecc.) e l'attributo \texttt{data-theme} sull'elemento
\texttt{<html>}, evitando la duplicazione di stili e rendendo i cambi di tema
istantanei senza ricaricare la pagina.

\subsubsection{Accessibilità WCAG}

L'interfaccia è progettata per essere utilizzabile da tutti:

\begin{itemize}
    \item \textbf{Attributi ARIA}: i dialog (modal login, modal recensione) usano
    \texttt{role="dialog"}, \texttt{aria-modal="true"} e \texttt{aria-labelledby}
    per una corretta lettura da parte degli screen reader.

    \item \textbf{Navigazione da tastiera}: le stelle di valutazione hanno
    \texttt{role="button"}, \texttt{tabindex="0"} e gestiscono
    \texttt{keydown} (Enter/Spazio), permettendo il voto senza mouse.

    \item \textbf{Skip-to-content}: un link nascosto ``Salta al contenuto''
    consente agli utenti di tastiera di bypassare la navbar. Utilizza
    \texttt{:focus-visible} invece di \texttt{:focus} per non apparire
    al tocco su dispositivi mobili.

    \item \textbf{Riduzione del movimento}: la media query
    \texttt{prefers-reduced-motion: reduce} disabilita tutte le animazioni
    CSS e JavaScript per gli utenti che ne fanno richiesta.

    \item \textbf{Stili di stampa}: la media query \texttt{@media print}
    nasconde navbar, bottoni e filtri, producendo ricevute stampabili pulite.
\end{itemize}

\subsubsection{Responsività mobile-first}

Il CSS è scritto con approccio \textbf{mobile-first}: gli stili base sono
ottimizzati per schermi piccoli, con override progressivi al crescere della
risoluzione. Su mobile compare una bottom navigation bar fissa con link
rapidi a Home, Aggiungi parcheggio e Profilo. Le griglie di card utilizzano
\texttt{CSS Grid} con \texttt{auto-fit} e \texttt{minmax()} per adattarsi
fluidamente a qualsiasi larghezza senza breakpoint rigidi.

\subsection{Sicurezza}

\subsubsection{Gestione delle credenziali}

Le credenziali del database non sono mai presenti nel codice sorgente. Un
parser leggero incluso in \texttt{config.php} legge le variabili da un
file \texttt{.env} escluso dal controllo di versione tramite \texttt{.gitignore}:

\begin{lstlisting}[language=PHP, caption={Caricamento sicuro delle variabili d'ambiente.}]
$envFile = dirname(__DIR__) . '/.env';
if (file_exists($envFile)) {
    foreach (file($envFile, FILE_IGNORE_NEW_LINES) as $line) {
        if (str_starts_with(trim($line), '#')) continue;
        [$k, $v] = array_map('trim', explode('=', $line, 2));
        if (!empty($k)) putenv("$k=$v");
    }
}
define('DB_PASS', getenv('DB_PASS') ?: '');
\end{lstlisting}

\subsubsection{Protezione contro attacchi comuni}

\begin{itemize}
    \item \textbf{SQL Injection}: ogni query al database utilizza prepared
    statements MySQLi con \texttt{bind\_param}, rendendo impossibile
    l'iniezione di codice SQL attraverso gli input utente.

    \item \textbf{Password sicure}: le password vengono cifrate con
    \texttt{password\_hash()} (algoritmo bcrypt, fattore di costo 12)
    e verificate con \texttt{password\_verify()}.

    \item \textbf{Rate limiting}: l'endpoint di login traccia i tentativi
    falliti nella sessione PHP e blocca ulteriori tentativi per 5 minuti
    dopo 5 errori consecutivi.

    \item \textbf{Validazione upload}: i file immagine caricati vengono
    verificati tramite la classe \texttt{finfo} (rilevamento MIME reale,
    non basato sull'estensione del file) e limitati a 5~MB.

    \item \textbf{XSS}: tutti i dati provenienti dal database vengono
    passati attraverso \texttt{htmlspecialchars()} prima di essere
    inseriti nel DOM lato PHP; lato JavaScript si usano le proprietà
    \texttt{textContent} per i valori dinamici.
\end{itemize}

\subsection{Generazione PDF}

Le ricevute di prenotazione vengono generate lato server con la libreria
open-source \textbf{dompdf}, che converte HTML/CSS in PDF senza dipendere
da servizi cloud. La ricevuta include: dati del cliente, dati del
proprietario del parcheggio, dettagli della prenotazione (data, orario,
importo totale) e badge dello stato. La funzione
\texttt{date\_default\_timezone\_set('Europe/Rome')} in \texttt{config.php}
garantisce che il timestamp di generazione rifletta il fuso orario italiano
indipendentemente dalla configurazione del server.

\section{Risultati}
\label{sec:risultati}

\subsection{Funzionalità implementate}

L'applicazione implementa un ciclo completo di utilizzo, dal primo accesso
alla generazione della ricevuta:

\begin{enumerate}[leftmargin=*, label=\textbf{\arabic*.}]
    \item \textbf{Registrazione e autenticazione}: form di registrazione con
    validazione lato client (formato email, lunghezza password minima 8 caratteri)
    e lato server (unicità email, cifratura bcrypt). Il login è protetto da
    rate limiting (max 5 tentativi per 5 minuti).

    \item \textbf{Home con filtri avanzati}: griglia di parcheggi con filtri
    combinabili in tempo reale (testo libero, prezzo massimo, valutazione minima,
    solo aperti oggi, ordinamento). I filtri operano sul DOM senza chiamate
    server aggiuntive. Il pulsante ``Carica altri'' carica i parcheggi
    successivi tramite AJAX mantenendo i filtri attivi.

    \item \textbf{Ricerca geografica}: il pulsante ``Cerca vicino a me'' ottiene
    le coordinate GPS tramite \texttt{navigator.geolocation} e interroga l'API
    con la formula di Haversine, restituendo i parcheggi entro un raggio
    configurabile (default 5~km).

    \item \textbf{Dettaglio parcheggio}: pagina con mappa interattiva Leaflet,
    orari, disponibilità per giorno della settimana, form di prenotazione con
    calcolo automatico dell'importo, sezione recensioni con media e voti.

    \item \textbf{Prenotazione}: validazione del giorno selezionato rispetto
    alla disponibilità del parcheggio sia lato client che server; controllo
    della disponibilità oraria (nessuna sovrapposizione con prenotazioni
    esistenti); transizione automatica di stato
    \texttt{in\_attesa} $\to$ \texttt{confermata} $\to$ \texttt{completata}
    (quest'ultima attivata al primo accesso al profilo dopo la data di scadenza).

    \item \textbf{Recensioni}: gli utenti possono lasciare una recensione
    (1--5 stelle + commento opzionale) per ogni prenotazione completata,
    sia dalla pagina del proprio profilo sia direttamente dalla pagina
    del parcheggio. È consentita una sola recensione per prenotazione
    (vincolo unico in database).

    \item \textbf{Profilo utente}: riepilogo prenotazioni con badge di stato,
    azioni per cancellare una prenotazione attiva, download della ricevuta PDF,
    gestione dei propri parcheggi pubblicati con possibilità di eliminazione.

    \item \textbf{Aggiunta parcheggio}: form con compilazione automatica
    dell'indirizzo via GPS (Nominatim), selezione dei giorni disponibili,
    orari di apertura/chiusura, upload foto con validazione MIME e
    dimensione massima (5~MB).

    \item \textbf{Ricevuta PDF}: documento generato lato server con dompdf
    contenente dati del cliente, dati del proprietario, dettaglio della
    prenotazione e importo totale. Include stili di stampa ottimizzati
    \texttt{(@media print)}.
\end{enumerate}

\subsection{Interfacce utente}

\subsubsection{Pagina Home}

La home presenta la griglia di parcheggi con:
\begin{itemize}
    \item Barra di ricerca testuale in navbar (funziona su tutte le pagine con
    redirect automatico alla home se non già presenti).
    \item Pannello filtri avanzati (collassabile, con pulsante toggle e attributo
    \texttt{aria-expanded} aggiornato dinamicamente).
    \item Badge ``Aperto oggi'' calcolato in PHP sulla base del giorno corrente.
    \item Contatore ``Mostrati X di Y parcheggi'' aggiornato dopo ogni ``Carica altri''.
\end{itemize}

\subsubsection{Tema chiaro e scuro}

\begin{table}[H]
\centering
\caption{Variabili CSS principali nei due temi.}
\renewcommand{\arraystretch}{1.3}
\begin{tabular}{lll}
\toprule
\textbf{Variabile} & \textbf{Tema chiaro} & \textbf{Tema scuro} \\
\midrule
\texttt{-{}-bg}      & \texttt{\#ffffff} & \texttt{\#0f172a} \\
\texttt{-{}-surface} & \texttt{\#f7f7f7} & \texttt{\#1e293b} \\
\texttt{-{}-accent}  & \texttt{\#185FA5} & \texttt{\#3b82f6} \\
\texttt{-{}-text}    & \texttt{\#1a1a1a} & \texttt{\#f1f5f9} \\
\texttt{-{}-muted}   & \texttt{\#6b6b6b} & \texttt{\#94a3b8} \\
\texttt{-{}-border}  & \texttt{rgba(0,0,0,.1)} & \texttt{rgba(255,255,255,.1)} \\
\bottomrule
\end{tabular}
\end{table}

\subsection{Metriche del progetto}

\begin{table}[H]
\centering
\caption{Dimensione dei componenti principali del progetto.}
\renewcommand{\arraystretch}{1.3}
\begin{tabular}{lrr}
\toprule
\textbf{File} & \textbf{Righe} & \textbf{Note} \\
\midrule
\texttt{assets/js/main.js}    & $\sim$780 & JS client-side (no dipendenze) \\
\texttt{assets/css/style.css} & $\sim$1{,}450 & CSS con variabili, dark mode, print \\
\texttt{classes/*.php}        & $\sim$800 & 6 classi PHP con PHPDoc \\
\texttt{api/*.php}            & $\sim$450 & 6 endpoint RESTful \\
\texttt{pages/*.php}          & $\sim$500 & 5 viste PHP \\
\midrule
\textbf{Totale}               & $\sim$3{,}980 & \\
\bottomrule
\end{tabular}
\end{table}

\section{Conclusioni}
\label{sec:conclusioni}

\subsection{Obiettivi raggiunti}

Il progetto ParkShare ha raggiunto tutti gli obiettivi prefissati:

\begin{itemize}
    \item È stata realizzata un'applicazione web completa che copre l'intero
    ciclo di utilizzo: dalla scoperta di un parcheggio alla prenotazione,
    dalla recensione alla ricevuta PDF.

    \item L'architettura RESTful è implementata con rispetto dei principi
    fondamentali (metodi HTTP semantici, statelessness, formato JSON uniforme).

    \item La componente cartografica è completamente open-source e gratuita,
    senza dipendenza da Google Maps o altri servizi commerciali.

    \item L'interfaccia è accessibile, responsiva e usabile su qualsiasi
    dispositivo e con qualsiasi preferenza di sistema (tema chiaro/scuro,
    riduzione del movimento).

    \item Il codice adotta pratiche di sicurezza consolidate:
    prepared statements, bcrypt, rate limiting, gestione sicura delle
    credenziali tramite \texttt{.env}.
\end{itemize}

\subsection{Considerazioni critiche}

Alcune limitazioni del progetto attuale meritano di essere riconosciute:

\begin{itemize}
    \item \textbf{Pagamenti}: il sistema attuale non integra un gateway di
    pagamento reale. Le prenotazioni vengono confermate automaticamente
    senza transazione economica, il che nella realtà richiederebbe
    l'integrazione con un servizio come Stripe o PayPal.

    \item \textbf{CSRF}: la protezione contro il Cross-Site Request Forgery
    non è stata implementata nella versione corrente. In produzione,
    ogni endpoint che modifica dati dovrebbe validare un token CSRF.

    \item \textbf{Notifiche}: non è previsto un sistema di notifica
    (email o push) per eventi come la conferma di una prenotazione o
    la ricezione di una recensione.
\end{itemize}

\subsection{Sviluppi futuri}

Gli sviluppi più interessanti per una versione futura del progetto sarebbero:

\begin{enumerate}
    \item \textbf{Integrazione pagamenti}: connessione con Stripe Checkout
    per gestire transazioni reali, con ricevuta fiscale automatica.

    \item \textbf{Notifiche email}: invio di email transazionali (conferma
    prenotazione, promemoria) tramite SMTP, senza dipendere da servizi
    a pagamento grazie a librerie come PHPMailer con server SMTP gratuiti.

    \item \textbf{Progressive Web App (PWA)}: aggiunta di un Service Worker
    per permettere l'installazione dell'applicazione come app nativa e
    il funzionamento offline di base.

    \item \textbf{Pannello statistiche proprietario}: dashboard per i
    proprietari dei parcheggi con grafici di prenotazioni, guadagni e
    valutazioni nel tempo.

    \item \textbf{CSRF protection}: implementazione di token sincronizzati
    per la protezione completa contro attacchi Cross-Site Request Forgery.
\end{enumerate}

\subsection{Riflessioni finali}

Lo sviluppo di ParkShare ha rappresentato un'occasione per applicare
concretamente le tecnologie studiate nel corso, affrontando problemi
reali: come bilanciare usabilità e prestazioni, come strutturare
un'API RESTful coerente, come proteggere un'applicazione dagli attacchi
più comuni, come garantire l'accessibilità senza sacrificare l'estetica.

La scelta consapevole di evitare framework e librerie pesanti — sia lato
client che lato server — ha reso il progetto più didatticamente
significativo: ogni meccanismo è implementato in modo esplicito e comprensibile,
senza ``magia nera'' nascosta nelle dipendenze.

L'adozione di tecnologie open-source per la componente cartografica
(Leaflet + OpenStreetMap) dimostra come sia possibile realizzare
applicazioni professionali senza costi di licensing e senza compromettere
la privacy degli utenti — un principio che diventa sempre più rilevante
nel panorama digitale contemporaneo.

\end{document}